%% kapa: exam=1 round=29 year=2018 period=1 form=A subject=economics subject_ko=경제학원론 seq=28 number=68
%% source: Q-Net 감정평가사 기출문제 (cst00309) artlSeq=5202808 file=extracted/1008/5202808/2018년 제29회 감정평가사 1차시험 기출문제/2018년 제29회 감정평가사 1차시험 기출문제/1교시 A형.pdf
%% conversion: kordoc 4.24.51 -> tools/gen1.py (auto); review=proofread-sonnet; flags=none
\begin{question}{68}
개방경제의 국민소득 결정모형이 아래와 같다. 정부지출( $G$)과 조세( $T$)를 똑같이 200에서 300으로 늘리면 균형국민소득은 얼마나 늘어나는가? (단, $Y$는 국민소득이다.)
\begin{framedtext}
소비함수: $C=300+0.6(Y-T)$ \par
투자함수: $I=200$ \par
정부지출: $G=200$ \par
조 세: $T=200$ \par
수 출: $EX=400$ \par
수 입: $IM=100+0.1(Y-T)$ \par
\end{framedtext}
\begin{choices}
  \item 0
  \item 50
  \item 100
  \item 200
  \item 250
\end{choices}
\end{question}
\answer{3}
